%%%%%%%%%% source scan idx 8 | folio 1 %%%%%%%%%%

\begin{center}
\underline{正则曲面上的相交}

\medskip
\underline{P. Deligne 著}
\end{center}

本讲第一部分先介绍相交理论的一些预备知识（见 $[5]\cup[7]$），然后给出 Raynaud [9] 的一个定理的特殊情形的 Artin 证明（1.13），并先验地证明一个 Néron [8] 曾逐一情形验证的定理（1.15）。[2] 中、从而第 IX 讲中用到的 [9] 的各定理，都较容易从（1.13）推出。

第二部分沿着同一思路，说明“意大利学派”[1]、[3] 如何计算曲面上的相交数。主要结果 2.9 和 2.11 也见于 Northcott [12]、[13]。最后，在 2.27、2.30 和 2.34 中，我提出一个猜想关系，将约化曲线奇点的完备（versal）形式模概形的维数、该奇点的不变量 $\delta$ 与衡量其对称性的一个不变量联系起来。

本讲的内容在讨论班后续各讲中不会用到。

\medskip
1. \underline{算术曲面上的相交数}

1.0. 本节始终设 $X$ 是一个二维正则概形（regular scheme），$Y$ 是 $X$ 上的约化除子（reduced divisor），$k$ 是 $H^0(Y,\underline{O}_Y)$ 的一个子域。假定 $Y$ 是 $k$ 上的固有曲线（EGA II 7.4.1）。下面给出这种情形的三个例子。

\underline{例} 1.1. 设 $(S,\eta,s)$ 是离散赋值环谱（trait），$f:X\longrightarrow S$ 是正则、固有且平坦的 $S$-概形，各纤维均为纯一维。取 $Y=(f^{-1}(s))_{\mathrm{red}}$ 及 $k=k(s)$。从 1.10 起，我们将只考虑这种情形。

%%%%%%%%%% source scan idx 9 | folio 2 %%%%%%%%%%

\underline{例} 1.2. 设 $X_o$ 是二维正规局部环的谱，$s$ 是 $X_o$ 的闭点，而 $f:X\longrightarrow X_o$ 是正则连通的 $X_o$-概形，在 $X_o$ 上固有，与 $X_o$ 不同，但在 $X_o\setminus\{s\}$ 上与 $X_o$ 一致。取 $Y=(f^{-1}(s))_{\mathrm{red}}$ 及 $k=k(s)$。

\underline{例} 1.3. 取 $X$ 为域 $k$ 上的完备正则曲面，$Y$ 为 $X$ 上的约化曲线。

回顾下述引理（SGA 6 IV，Ex.~1.9）。

\underline{引理} 1.4. \underline{设 $Z_o$ 是 Noether 概形 $Z$ 的闭子概形。则 $Z_o$ 到 $Z$ 的包含，诱导从凝聚 $\underline{O}_{Z_o}$-模范畴的 Grothendieck 群到 $Z$ 上支集包含于 $Z_o$ 的凝聚层范畴的 Grothendieck 群的同构。}

由此引理，可以对 $X$ 上支集包含于 $Y$ 的凝聚层 $\mathcal F$ 定义 Euler--Poincaré 特征数 $\chi(\mathcal F)$，使得：

(i) 若 $\mathcal F$ 是凝聚 $\underline{O}_Y$-模层，则
\begin{equation}\tag{1.4.1}
\chi(\mathcal F)=\dim_k H^0(\mathcal F)-\dim_k H^1(\mathcal F).
\end{equation}

(ii) $\chi$ 对短正合列具有可加性。

若 $\underline{O}_X$-模复形 $K$ 的各上同调层都是支集包含于 $Y$ 的凝聚层，则置
\begin{equation}\tag{1.4.2}
\chi(K)=\sum_i(-1)^i\chi(\underline{H}^i(K)).
\end{equation}
若 $\mathcal F$ 的支集包含于 $Y$ 的一点 $y$，则
\begin{equation}\tag{1.4.3}
\chi(\mathcal F)=\operatorname{length}_{\underline{O}_y}(\mathcal F_y)[k(y):k].
\end{equation}

%%%%%%%%%% source scan idx 10 | folio 3 %%%%%%%%%%

在（1.1）、（1.2）或（1.3）的假设下，有
\begin{equation}\tag{1.4.4}
\chi(\mathcal F)=\operatorname{length}f_*\mathcal F-
\operatorname{length}R^1f_*\mathcal F.
\end{equation}
设 $D$ 是支集包含于 $Y$ 的除子 $\geq0$，$\mathcal L$ 是 $D$ 上的可逆层。则可用唯一的方式定义次数 $\deg_D(\mathcal L)$，使得
\begin{equation}\tag{1.4.5}
\left\{
\begin{array}{l}
\deg_D(\mathcal L'\otimes\mathcal L'')=
\deg_D(\mathcal L')+\deg_D(\mathcal L''),\\[3ex]
\text{对 }\mathcal L\subset\underline O,
\text{有 }\deg_D(\mathcal L)=-\chi(\underline O/\mathcal L).
\end{array}
\right.
\end{equation}
于是 $D$ 上的 Riemann--Roch 公式为
\begin{equation}\tag{1.4.6}
\chi(\mathcal L)=\chi(\underline{O}_D)+\deg_D(\mathcal L).
\end{equation}
若 $D=\sum n_iD_i$ 是 $D$ 分解为约化不可约除子的表达式，则
\begin{equation}\tag{1.4.7}
\deg_D(\mathcal L)=\sum n_i\deg_{D_i}(\mathcal L).
\end{equation}

\underline{定义} 1.5. \underline{设 $D$ 与 $E$ 是 $X$ 上的两个除子 $\geq0$，并假设 $D$ 的支集包含于 $Y$。$D$ 与 $E$ 的相交数（相对于域 $k$（1.0））定义为}
\begin{equation*}
(D,E)=\chi(\underline{O}_D\overset{\mathbb L}{\otimes}\underline{O}_E)
=\chi(\underline{O}_D\otimes\underline{O}_E)
-\chi(\operatorname{Tor}^1(\underline{O}_D,\underline{O}_E)).
\end{equation*}
事实上，在这里的假设下，$\operatorname{Tor}^2(\underline{O}_D,\underline{O}_E)=0$。

\underline{命题} 1.6. \underline{在 1.5 的假设下，有：}

(i) \quad \underline{相交数 $(D,E)$ 对除子 $D$ 与 $E$ 均具有可加性；借助线性性，可据此定义两个未必为正的除子的相交数。}

%%%%%%%%%% source scan idx 11 | folio 4 %%%%%%%%%%

(ii) \quad $(D,E)=\deg_D(\underline O(E))$.

(iii) \quad \underline{若 $D$ 与 $E$ 仅在有限多个点相交，则 $(D,E)=\chi(\underline{O}_D\otimes\underline{O}_E)\geq0$，且仅当 $D$ 与 $E$ 不交时等号才成立。}

(iv) \quad \underline{若 $D$ 与 $E$ 的支集均包含于 $Y$，则 $(D,E)=(E,D)$。}

(v) \quad \underline{若 $N$ 是 $D$ 的余法层（conormal sheaf），则}
\begin{equation*}
D^2=(D,D)=-\deg_D(N).
\end{equation*}

(vi) \quad \underline{若未必为正的 $E$ 是 $X$ 上某个亚纯函数的除子，则 $(D,E)=0$。}

\underline{证明}。正合列
\begin{equation*}
0\longrightarrow\underline O(-E)\longrightarrow\underline O_X
\longrightarrow\underline O_E\longrightarrow0
\end{equation*}
结合（1.4.6）给出
\begin{equation*}
\begin{array}{rl}
(D,E)\overset{\mathrm{def}}{=}&
\chi(\underline{O}_D\overset{\mathbb L}{\otimes}\underline{O}_E)\\
=&\chi(\underline{O}_D\overset{\mathbb L}{\otimes}\underline{O}_X)
-\chi(\underline{O}_D\overset{\mathbb L}{\otimes}\underline O(-E))\\[1ex]
=&\chi(\underline{O}_D)-\chi(\underline O(-E)|D)
=-\deg_D(\underline O(-E))=\deg_D(\underline O(E)).
\end{array}
\end{equation*}
由（1.4.5）和（1.4.7），这就证明了 (ii)、(vi) 及双可加性。在 (iii) 的假设下，$\operatorname{Tor}^1(\underline{O}_D,\underline{O}_E)=0$，因而得到结论。在 (v) 的假设下，有
\begin{equation*}
\underline{O}_D\otimes\underline{O}_D=\underline{O}_D
\quad\text{且}\quad
\operatorname{Tor}^1(\underline{O}_D,\underline{O}_D)\simeq N
\quad(\text{SGA 6 VII 2.5}),
\end{equation*}
所以
\begin{equation*}
(D,D)=\chi(\underline{O}_D)-\chi(N)=-\deg_D(N).
\end{equation*}
最后，(iv) 由张量积的对称性得到。

注意，即使不假设除子为正，(ii)、(iv) 和 (vi) 仍然成立。

%%%%%%%%%% source scan idx 12 | folio 5 %%%%%%%%%%

\underline{命题} 1.7. \underline{设 $(Y_i)_{1\leq i\leq k}$ 是 $Y$ 的各不可约分支，$f$ 是 $X$ 上的亚纯函数，其除子为}
\begin{equation*}
\operatorname{div}(f)=\sum_i n_iY_i+Z
\end{equation*}
\underline{其中 $|Z|\cap Y$ 为有限集。置 $D_i=n_iY_i$，则对 $(a_i)\in\mathbb Q^k$ 有：}
\begin{equation*}
\left(\sum_i a_iD_i\right)^2
=-\sum_{i<j}(a_i-a_j)^2(D_i,D_j)-\sum_i a_i^2(D_i,Z).
\end{equation*}

\underline{证明}（依照 [7]）。由 1.6 (v)、(vi)，
\begin{equation*}
\begin{aligned}
\left(\sum_i a_iD_i\right)^2
&=\sum_i a_i\left(D_i,\sum_j a_jD_j\right)
=\sum_i a_i\left(D_i,\sum_j a_jD_j-a_i\operatorname{div}f\right)\\
&=\sum_i a_i\left(D_i,\sum_{j\ne i}(a_j-a_i)D_j-a_iZ\right)\\
&=\sum_{i\ne j}a_i(a_j-a_i)(D_i,D_j)-\sum_i a_i^2(D_i,Z)\\
&=-\sum_{i<j}(a_i-a_j)^2(D_i,D_j)-\sum_i a_i^2(D_i,Z).
\end{aligned}
\end{equation*}

\underline{推论} 1.8. \underline{在例 1.1 的假设下，取 $f$ 为 $S$ 的一个一致化参数 $\pi$。此时 $Z=0$ 且}
\begin{equation*}
\left(\sum_i k_iY_i\right)^2
=-\sum_{i<j}\frac{1}{n_in_j}(k_in_j-k_jn_i)^2(Y_i,Y_j);
\end{equation*}
\underline{所以相交型为负半定；若 $Y$ 连通，则其核恰由 $\operatorname{div}(\pi)$ 的倍数组成。}

\underline{推论} 1.9（Mumford [7]）。\underline{在例 1.2 的假设下，取 $f$ 为局部环 $\underline O_s$ 的极大理想中的非零元。此时 $Z>0$，所以相交型为负定。}